\scanpage{34}
défini par l'aut \deleted{\unsure{de}} $\tau_X$ de $X(K)$. Ainsi, un motif sur $\mathbf R$ définit :

a) Une structure de Hodge $V$ sur $K$

b) Un autom $f_{\mathbf R}$ \inserted{de carré 1} du $K$-espace vectoriel \deleted{sousjacent} \inserted{$V_K$}, cet autom étant ``réel'' i.e. qui transforme $V_{\mathbf R}$ dans $V_{\mathbf R}$, et $\mathbf Q$-\uline{rationnel} i.e. transforme $V$ en $V$

\deleted{\unsure{La réciproque|La conjugaison}}

Ces deux données étant assujetties à la condition
\[
\left\{\begin{array}{l}
(f_{\mathbf R})\sigma\ \text{\unsure{invarie} la filtration naturelle}\\
\text{i.e. }f_{\mathbf R}\ \text{échange la filtration donnée et l'opposée}
\end{array}\right.
\]
En termes de ``\unsure{donnée de descente}'', on peut \deleted{\unsure{dire aussi}} \inserted{exprimer} aussi b) en disant qu'on a un isom de la structure de Hodge $V$ avec la structure de Hodge \deleted{\unsure{qui}} ``conjuguée'' : c'est le frobenius réel.

Soit $G$ le groupe de Galois motivique de $M$ sur $\mathbf R$. Alors $f_{\mathbf R}$ correspond à un élément
\[
f_{\mathbf R}\in G(\mathbf Q)
\]
appelé \uline{frobenius réel}. C'est un élément d'ordre 2. Son image dans $G/G^0$ engendre ce groupe. Cet élément normalise $\underline i:S_{\mathbf R}\longrightarrow G_{\mathbf R}$, de façon précise
\[
f_{\mathbf R}i(\lambda)f_{\mathbf R}^{-1}=i(\deleted{\unsure{\lambda}}\lambda')
\qquad\text{où }\lambda\mapsto\lambda'\text{ est la symétrie canonique de }S_{\mathbf R}.
\]
Donc sur $U_{\mathbf R}$ on a
\[
(*)\qquad f_{\mathbf R}i(\lambda)f_{\mathbf R}^{-1}=i(\lambda^{-1})
\qquad (\lambda\in U_{\mathbf R}\deleted{\unsure{(\mathbf R)}})
\]
Le foncteur $M\mapsto(\text{structure de Hodge associée, avec }f_{\mathbf R})$ est pleinement fidèle en vertu des conj. de Hodge.

\margin{\unsure{Les structures en jeu sont donc} $G$ sur $\mathbf Q$, $\varepsilon:G\to G_m$, $i:S_{\mathbf R}\unsure{-|\longrightarrow}G_{\mathbf R}$ (avec $i\vert G_{m\mathbf R}$ défini sur $\mathbf Q$ etc) $+f_{\mathbf R}\in G(\mathbf Q)$, satisfaisant les conditions \unsure{exp.}}
\mspage{33}
